% 图契约：树上的每条边都是唯一通路；图中的端口、几何、箭头和方程对应本题。
% 版面：沿用前三模块的几何图+推导分区；目标印宽17cm，简图至中等密度。
% 布局：左侧真实结构，右侧关键关系；真实电路边不作为架构流程箭头处理。
\documentclass[tikz,border=7pt]{standalone}
\usepackage{fontspec,xeCJK,amsmath,amssymb}
\IfFontExistsTF{Times New Roman}{\setmainfont{Times New Roman}}{\setmainfont{TeX Gyre Termes}}
\IfFontExistsTF{Heiti SC}{\setCJKmainfont[AutoFakeBold=1.5]{Heiti SC}}{\setCJKmainfont[AutoFakeBold=1.5]{FandolHei-Regular}}
\usetikzlibrary{arrows.meta,bending,calc,patterns,positioning,decorations.pathmorphing}
\definecolor{ink}{HTML}{243745}
\definecolor{blue}{HTML}{275DAB}
\definecolor{green}{HTML}{177C63}
\definecolor{red}{HTML}{BD4935}
\definecolor{orange}{HTML}{AA6C16}
\definecolor{purple}{HTML}{715393}
\tikzset{every node/.style={font=\fontsize{9.3}{12.5}\selectfont,text=ink},
arrow/.style={-{Stealth[length=5pt 1.5,width'=0pt 3]},shorten >=2pt,shorten <=1pt,line width=1pt},
arrow short/.style={-{Stealth[length=3pt 1.0,width'=0pt 2.5]},shorten >=1pt,shorten <=0pt,line width=.8pt},
axis/.style={arrow,draw=ink!60,line width=.6pt},
vec/.style={arrow,draw=blue},
wire/.style={draw=ink!75,line width=.8pt},
guide/.style={dashed,draw=ink!45,line width=.5pt},
card/.style={rounded corners=2pt,draw=blue!65,fill=blue!4,align=center,text width=6.5cm,inner xsep=9pt,inner ysep=7pt},
title/.style={font=\fontsize{10.5}{14}\selectfont\bfseries,text width=15cm,align=center},
label/.style={font=\fontsize{8.4}{11}\selectfont,align=center},
note/.style={label,text width=6.1cm}}
\begin{document}
\begin{tikzpicture}[x=1cm,y=1cm]
\path[use as bounding box] (-8,-3.5) rectangle(8,3.65);
\node[title] at (0,3.25) {M4-064\quad 树上的每条边都是唯一通路};
\begin{scope}[shift={(-4.05,0)}]
\draw[wire,draw=ink](-2.3,0)--(-1.8249999999999997,0.0) (-1.275,0.0)--(-0.8,0);\begin{scope}[shift={(-1.5499999999999998,0.0)},rotate=0.0]\draw[draw=ink,fill=white,line width=.8pt](-0.275,-.085)rectangle(0.275,.085);\end{scope}
\draw[wire,draw=ink](-0.8,0)--(-0.21343135003197305,0.4765870280990219) (0.21343135003197305,0.8234129719009782)--(0.8,1.3);\begin{scope}[shift={(0.0,0.65)},rotate=39.0938588862295]\draw[draw=ink,fill=white,line width=.8pt](-0.275,-.085)rectangle(0.275,.085);\end{scope}
\draw[wire,draw=ink](-0.8,0)--(-0.21343135003197305,-0.4765870280990219) (0.21343135003197305,-0.8234129719009782)--(0.8,-1.3);\begin{scope}[shift={(0.0,-0.65)},rotate=-39.0938588862295]\draw[draw=ink,fill=white,line width=.8pt](-0.275,-.085)rectangle(0.275,.085);\end{scope}
\draw[wire,draw=ink](0.8,1.3)--(1.275,1.3) (1.8249999999999997,1.3)--(2.3,1.3);\begin{scope}[shift={(1.5499999999999998,1.3)},rotate=0.0]\draw[draw=ink,fill=white,line width=.8pt](-0.275,-.085)rectangle(0.275,.085);\end{scope}
\draw[wire,draw=ink](0.8,-1.3)--(1.275,-1.3) (1.8249999999999997,-1.3)--(2.3,-1.3);\begin{scope}[shift={(1.5499999999999998,-1.3)},rotate=0.0]\draw[draw=ink,fill=white,line width=.8pt](-0.275,-.085)rectangle(0.275,.085);\end{scope}
\fill[ink] (-2.3,0)circle(2pt);\node[left=4pt] at(-2.3,0){$a$};
\fill[ink] (-0.8,0)circle(2pt);
\fill[ink] (0.8,1.3)circle(2pt);
\fill[ink] (0.8,-1.3)circle(2pt);
\fill[ink] (2.3,1.3)circle(2pt);
\fill[ink] (2.3,-1.3)circle(2pt);\node[below=4pt] at(2.3,-1.3){$b$};
\node[note] at (0,-2.75) {示意六节点树；结论允许每边阻值不同};
\end{scope}
\node[card,draw=blue!65,fill=blue!4] at (3.9,1.8) {删除任一边都会使树断开};
\node[card,draw=green!65,fill=green!4] at (3.9,0) {$R_e=r_e$};
\node[card,draw=orange!65,fill=orange!4] at (3.9,-1.8) {$\sum_eR_e/r_e=N-1$};
\end{tikzpicture}\end{document}